% Part: normal-modal-logic
% Chapter: sequent-calculus
% Section: rules-for-K

\documentclass[../../../include/open-logic-section]{subfiles}

\begin{document}
\olfileid{nml}{seq}{rul}

\olsection{\Log{K}的规则}

常规命题联结词的规则与标准相继式演算~$\Log{LK}$ 相同。公理也一样：任何形如 $!A \Sequent !A$ 的相继式都算作公理。

对于模态算子\iftag{prvBox}{\iftag{prvDiamond}{~$\Box$ 和}{~$\Box$}}{}\iftag{prvDiamond}{~$\Diamond$}{}，我们有以下附加\iftag{notprvBox,notprvDiamond}{规则}{规则}：
\[
  \iftag{prvBox}{\iftag{prvDiamond}{% <> and [] primitive
    \Axiom$\Gamma \fCenter \Delta, !A$
    \RightLabel{$\Box$}
    \UnaryInf$\Box\Gamma \fCenter \Diamond\Delta, \Box!A$
    \DisplayProof
    \qquad
    \Axiom$!A, \Gamma \fCenter \Delta$
    \RightLabel{$\Diamond$}
    \UnaryInf$\Diamond!A, \Box\Gamma \fCenter \Diamond\Delta$
    \DisplayProof
}{% only Box primitive
\Axiom$\Gamma \fCenter !A$
\RightLabel{$\Box$}
\UnaryInf$\Box\Gamma \fCenter \Box!A$
\DisplayProof
}}{% only <> primitive
  \Axiom$!A \fCenter \Delta$
  \RightLabel{$\Diamond$}
  \UnaryInf$\Diamond!A \fCenter \Diamond\Delta$
  \DisplayProof
}
\]
这里，\iftag{prvBox}{$\Box\Gamma$ 表示在~$\Gamma$ 中的每个公式前加上 $\Box$ 后得到的公式序列\iftag{prvDiamond}{，而}{}}{}\iftag{prvDiamond}{$\Diamond\Delta$ 表示在~$\Delta$ 中的每个公式前加上 $\Diamond$ 后得到的公式序列}{}。\iftag{notprvDiamond}{$\Box$ 规则的前提的右端必须至多包含一个公式~$!A$。}{}%
\iftag{notprvBox}{$\Diamond$ 规则的前提的左端必须至多包含一个公式~$!A$。}{}%
\iftag{prvBox}{$\Gamma$\iftag{prvDiamond}{~和~}{}}{}\iftag{prvDiamond}{$\Delta$}{}可以为空；此时，结论相继式中对应的部分
\iftag{prvBox}{$\Box\Gamma$\iftag{prvDiamond}{~和~}{}}{}\iftag{prvDiamond}{$\Diamond\Delta$}{}也同样为空。

\iftag{prvBox}{在右侧仅对单个公式~$!A$ 添加 $\Box$\iftag{prvDiamond}{，以及}{}}{}\iftag{prvDiamond}{在左侧仅对单个公式~$!A$ 添加 $\Diamond$}{}，这一限制是必要的。如果我们允许
\iftag{prvBox}{在右侧为任意多个公式添加 $\Box$\iftag{prvDiamond}{，或}{}}{}\iftag{prvDiamond}{在左侧为任意多个公式添加 $\Diamond$}{}，我们就能够推导出：
  \[
    \iftag{prvBox}{
      \Axiom$!A \fCenter !A$
      \RightLabel{\RightR{\lnot}}
      \UnaryInf$\fCenter !A, \lnot !A$
      \RightLabel{$\Box*$}
      \UnaryInf$\fCenter \Box!A, \Box\lnot !A$
      \doubleLine
      \RightLabel{\RightR{\lor}}
      \UnaryInf$\fCenter \Box!A \lor \Box \lnot !A$
      \DisplayProof}{}
    \iftag{notprvBox,notprvDiamond}{}{\qquad}
    \iftag{prvDiamond}{
      \Axiom$!A \fCenter !A$
      \RightLabel{\LeftR{\lnot}}
      \UnaryInf$\lnot !A, !A \fCenter$
      \RightLabel{$\Diamond*$}
      \UnaryInf$\Diamond\lnot !A,\Diamond!A \fCenter $
      \RightLabel{\RightR{\lnot}}
      \UnaryInf$\Diamond!A \fCenter \lnot\Diamond\lnot !A$
      \RightLabel{\RightR{\lif}}
      \UnaryInf$ \fCenter \Diamond!A \lif \lnot\Diamond\lnot !A$
      \DisplayProof}{}
  \]
但 \iftag{prvBox}{$\Box!A \lor \Box \lnot !A$\iftag{prvDiamond}{ 和 $\Diamond!A \lif \lnot\Diamond\lnot !A$ 在~$\Log{K}$ 中并非}{在~$\Log{K}$ 中并非}}{$\Diamond!A \lif \lnot\Diamond\lnot !A$ 在~$\Log{K}$ 中并非}有效的。

如果我们允许前提中除了~$!A$ 之外还有辅助公式，并且允许
\iftag{prvBox}{$\Box$ 规则仅在右侧为~$!A$ 添加 $\Box$\iftag{prvDiamond}{，或允许}{}}{}\iftag{prvDiamond}{$\Diamond$ 规则仅在左侧为~$!A$ 添加 $\Diamond$}{}（而对辅助公式不作处理），我们就能够推导出：
\[
  \iftag{prvBox}{
    \Axiom$!A \fCenter !A$
    \RightLabel{\RightR{\lnot}}
    \UnaryInf$\fCenter !A, \lnot !A$
    \RightLabel{\RightR{\Exchange}}
    \UnaryInf$\fCenter \lnot !A, !A$
    \RightLabel{$\Box*$}
    \UnaryInf$\fCenter \lnot!A, \Box !A$
    \doubleLine
    \RightLabel{\RightR{\lor}}
    \UnaryInf$\fCenter \lnot!A \lor \Box !A$
    \DisplayProof}{}
  \iftag{notprvBox,notprvDiamond}{}{\qquad}
  \iftag{prvDiamond}{
    \Axiom$!A \fCenter !A$
    \RightLabel{\LeftR{\lnot}}
    \UnaryInf$\lnot !A, !A \fCenter$
    \RightLabel{$\Diamond*$}
    \UnaryInf$\Diamond\lnot !A, !A \fCenter $
    \RightLabel{\RightR{\lnot}}
    \UnaryInf$!A \fCenter \lnot\Diamond\lnot !A$
    \RightLabel{\RightR{\lif}}
    \UnaryInf$ \fCenter !A \lif \lnot\Diamond\lnot !A$
    \DisplayProof}{}
\]
但 \iftag{prvBox}{$\lnot!A \lor \Box !A$（等价于 $!A \lif \Box!A$）\iftag{prvDiamond}{ 和 $!A \lif \lnot\Diamond\lnot !A$ 在~$\Log{K}$ 中并非}{在~$\Log{K}$ 中并非}}{$!A \lif \lnot\Diamond\lnot !A$ 在~$\Log{K}$ 中并非}有效的。

\end{document}